\chapter{Static Memory: Bounding the Stack and Heap}
\label{ch:static-mem}\index{stack}\index{memory!static bounding}

A workstation program runs in virtual memory: the stack grows until it hits a
guard page, the heap until the OS says no. A bare-metal ESP32-S3 has neither
luxury. There is a fixed amount of internal SRAM -- a few hundred kilobytes --
divided at link time into \texttt{.data}, \texttt{.bss}, the task stacks and
whatever is left for a heap, with the external PSRAM as optional overflow. Nothing
grows on demand; if a stack runs past its allotment it corrupts whatever is below
it. So a real question on this chip is: \emph{how much memory does this program
use, worst case, and does it fit?} This chapter is the accounting.

The three runtime profiles (Chapter~\ref{ch:profiles}) make very different
promises here, so the chapter keeps coming back to them: \texttt{light-tasking} is
the austere one -- no heap, no secondary stack, no finalization -- while
\texttt{embedded} and \texttt{full} trade that austerity for a real heap and full
exceptions.

\section{The environment-task stack}\index{environment task}

The first stack is the one \texttt{Main} runs on: the environment task's stack. Its
size is a build-time knob, \texttt{ENV\_STACK\_SIZE}, defaulting to 32\,KiB under
\texttt{embedded}/\texttt{full} and 16\,KiB under \texttt{light-tasking} (the
exception machinery needs more headroom, so the heavier profiles start larger). The
build turns the number into the linker symbols \texttt{\_\_stack\_start} and
\texttt{\_\_stack\_end} that \texttt{start.S} loads into the stack pointer; an
example that recurses deeply just exports a bigger value (the LISP and SIMD examples
set 64\,KiB and 128\,KiB).

\begin{lstlisting}[style=sh]
export ENV_STACK_SIZE=131072    # 128 KiB env-task stack for this example
\end{lstlisting}

When even that is not enough -- the full-profile ACATS sweep elaborates very deep --
the \texttt{ENV\_STACK\_PSRAM} option moves the whole env stack to the top of the
8\,MB PSRAM, trading a little access latency for a great deal of room. The
bootloader maps PSRAM \emph{before} \texttt{start.S} sets the stack pointer
(Chapter~\ref{ch:bootloader}), so this is safe at the very first instruction.

\section{Per-task stacks}\index{task!stack size}

Every other task gets its own fixed stack, sized by the \texttt{Storage\_Size}
aspect on the task or, if you say nothing, the runtime default from
\texttt{System.Parameters} -- 20\,KiB under \texttt{light-tasking}, 12\,KiB under
\texttt{full} (smaller there only so the ACATS sweep can hold more tasks at once).
A task that knows it is shallow should say so and reclaim the rest:

\begin{lstlisting}
task type Poller with Storage_Size => 4 * 1024;   -- 4 KiB is plenty
\end{lstlisting}

Because the Jorvik/Ravenscar model (Chapter~\ref{ch:tasking}) forbids dynamic task
creation, the \emph{set} of tasks is known at compile time, so the sum of all task
stacks is a static quantity -- you can add it up. There is no pool of tasks coming
and going to reason about; what you declare is what you pay.

\section{The secondary stack}\index{secondary stack}

Ada returns unconstrained results -- a \texttt{String} of computed length, say --
on a \emph{secondary} stack, separate from the primary one. Its per-task size is set
by the binder \texttt{-D} switch you will see in every project file
(\texttt{-D8k}, \texttt{-D16k}); the default is a modest 4\,KiB.

\begin{lstlisting}[style=sh]
for Switches ("Ada") use ("-D16k", "-Q2", "-Mada_main");
\end{lstlisting}

The neighbouring \texttt{-Q} switch is subtler and worth knowing. Under
\texttt{-Q2} (the bareboard default) each task's secondary stack is handed out from
a \emph{monotonic pool} that is never reclaimed -- fine for the steady-state task
set, but a program that created and destroyed tasks forever would exhaust it. The
\texttt{full} profile carries a runtime patch that makes secondary stacks
heap-managed instead (built \texttt{-Q0}), so they are freed with the task; it is
one of the concrete differences between the profiles, not just a flag.

\section{The heap --- or the absence of one}\index{heap}\index{TLSF}

This is where the profiles split hardest. \texttt{light-tasking} has \emph{no heap}:
it sets \texttt{No\_Allocators}, and its \texttt{s-memory} ``allocator'' is a bump
pointer whose \texttt{Free} raises \texttt{Program\_Error}. Memory is whatever you
declared statically; \texttt{new} is a compile-time error. That is a feature -- it
is the profile you reach for when you want a provably bounded memory footprint and
no fragmentation, ever.

\texttt{embedded} and \texttt{full} provide a real heap: a TLSF (Two-Level
Segregated Fit) allocator with O(1) \texttt{malloc}/\texttt{free} and bounded
fragmentation, described in full in Chapter~\ref{ch:heap}. Its arena is placed by
the build -- by default the leftover internal DRAM after \texttt{.data},
\texttt{.bss} and the stacks (a couple of hundred KiB), or, with
\texttt{HEAP\_PSRAM=1}, the 8\,MB PSRAM for programs that need real room (the LISP
arena and the ext4 block cache both live there).

\begin{lstlisting}[style=sh]
export HEAP_SIZE=1 HEAP_PSRAM=1   # heap arena in PSRAM, not internal DRAM
\end{lstlisting}

\section{What the profile forbids}

Memory bounding is really enforced by \emph{restrictions} -- pragmas the profile's
\texttt{system.ads} sets, which make the expensive constructs compile errors rather
than runtime surprises. The contrast is the whole story:

\begin{center}
\begin{tabular}{l c c}
\toprule
\textbf{Restriction} & \textbf{light-tasking} & \textbf{embedded / full} \\ \midrule
\texttt{No\_Allocators} (no heap)            & yes & no \\
\texttt{No\_Secondary\_Stack}                & yes & no \\
\texttt{No\_Finalization}                    & yes & no \\
\texttt{No\_Exception\_Propagation}          & yes & no \\
\texttt{No\_Implicit\_Dynamic\_Code}         & yes & no\,$^{*}$ \\
\bottomrule
\end{tabular}
\end{center}

\noindent$^{*}$\,\texttt{embedded} relaxes it but still expects closure-free,
library-level callbacks (Chapter~\ref{ch:driver-design}); \texttt{full} re-points
nested-subprogram trampolines into executable memory at task creation instead
(Chapter~\ref{ch:runtime-build}).

Under \texttt{light-tasking} the combination means a program's entire memory use --
every stack, no heap, no hidden secondary-stack growth -- is fixed at link time and
checkable by the compiler. Each step up to \texttt{embedded} and \texttt{full} buys
expressiveness (exceptions you can catch, controlled types, dynamic data) at the
cost of a quantity you now have to \emph{budget} rather than have proven zero.

\section{Doing the accounting}\index{stack!analysis}\index{tooling!x stack}

The static quantities above -- the task set, each stack size, the heap-arena
bounds, the restriction set -- are all pinned at link time, so they can be
\emph{measured} rather than guessed. The repository wires three tools onto the
\texttt{x} command for exactly that.

\subsection{Per-frame and worst-case: \texttt{x stack}}

\texttt{AdaCore}'s \texttt{GNATstack} is a commercial tool and is not in the
free Xtensa toolchain, but the underlying GCC switches are: \texttt{-fstack-usage}
emits one \texttt{.su} file per unit giving every subprogram's frame size and
whether it is \texttt{static}, \texttt{dynamic} or \texttt{bounded}, and
\texttt{-fcallgraph-info} emits the call graph with those frame sizes attached.
The build turns them on only when \texttt{STACK\_ANALYSIS=1} (a normal build stays
byte-identical), and

\begin{lstlisting}[language=bash]
./x stack <example> [--top N]
\end{lstlisting}

\noindent
forces an analysis rebuild and reports two things. First, the biggest single
frames, with any \texttt{dynamic}/\texttt{bounded} ones (a variable-length array or
\texttt{alloca} whose size is not a compile-time constant) called out -- those are
the frames that defeat static bounding. Second, the deepest stack-summed call chain
from each entry root: because Ravenscar/Jorvik forbids recursion the application's
call graph is a DAG, so the longest path is a genuine bound on the code GCC
compiled. The figure covers application units only -- the pinned runtime is built
without the switches, so calls into it are marked \texttt{+ext} and excluded; the
watermark below closes that gap. If the analyser \emph{does} find a cycle it says
so loudly (\texttt{!!RECURSIVE}) and downgrades that chain's number to a lower
bound, because a recursive depth cannot be bounded statically.

\subsection{The measured peak: stack painting}\index{stack!high-water mark}\index{stack!painting}

Static analysis cannot see the runtime, the C startup, the interrupt path or
hand-written assembly. To account for those you measure the real peak by
\emph{painting}: fill the unused stack with a sentinel word, run a representative
workload, then scan for the lowest word the program actually overwrote. The gap
from there to the top of the stack is the high-water mark. \texttt{ESP32S3.Stack\_Usage}
does this for the environment task in a few lines --

\begin{lstlisting}[language=ada]
ESP32S3.Stack_Usage.Paint_Env_Stack;   --  once, as early in Main as possible
--  ... run the workload ...
ESP32S3.Stack_Usage.Report;            --  "stack: env used=5184 free=11200 total=16384 (31%)"
\end{lstlisting}

\noindent
using the linker's \texttt{\_\_stack\_start}/\texttt{\_\_stack\_end} symbols as the
region bounds; the same \texttt{Paint}/\texttt{High\_Water} primitives take explicit
bounds so a library-level task can measure its own stack the same way. Because a
sentinel-valued word the program legitimately wrote looks pristine, the scan stops
at the first overwritten word from the bottom and so never \emph{under}-reports the
peak -- read it as ``at least this much.'' The \texttt{esp32s3\_stack\_usage}
example demonstrates the jump as a deep call runs, and \texttt{./x stack <example>
-{}-run} flashes it and captures the \texttt{stack:} line over the serial console,
so the static estimate and the measured peak sit side by side.

\subsection{Where the rest of the memory went: \texttt{x mem}}\index{tooling!x mem}

Stacks are one line item; \texttt{./x mem <example>} prints the rest -- the loaded
section sizes from the linked ELF, split into flash (code + read-only data), the
DRAM \texttt{.data}/\texttt{.bss}, the IRAM image and any reserved PSRAM, alongside
the configured \texttt{board.ads} bounds. Between \texttt{x stack} and \texttt{x mem}
the whole static budget -- every byte of flash and RAM, and the worst-case and
measured depth of every stack -- is visible from two commands. Turning those numbers
into a signed-off worst-case for a hard-real-time deployment is still your
judgement call (headroom, the \texttt{+ext} runtime contribution, the margin on a
\texttt{dynamic} frame), but the numbers themselves are now produced for you, not
reasoned about on paper.

\subsection{A worked example: the HTTPS client}

The \texttt{esp32s3\_tls\_weather} example -- a pure-Ada TLS~1.3 client fetching a
live forecast over the W5500 Ethernet stack -- is the most demanding example in the
tree, so it is the one worth putting under the tools. \texttt{./x stack tls\_weather}
reports:

\begin{lstlisting}
== Per-frame stack usage (application units) ==
   13 frames; total of all frames = 5 344 B (NOT a depth)
      bytes  qualifier        location / subprogram
      4 832  dynamic          main.adb:55  Main
        128  static           w5500_dev.adb:13  Bring_Up
         48  static           b__main.adb:281  ada_main
        ...
   !! 1 DYNAMIC/BOUNDED frame(s) -- static depth is NOT guaranteed:
      4 832  dynamic          main.adb:55  Main

== Worst-case call-chain depth per entry root ==
   5 008 B +ext   ada_main
          ada_main -> ada_main_program -> Bring_Up
\end{lstlisting}

\noindent
One number dominates everything: \texttt{Main} carries a \textbf{4\,832-byte,
\texttt{dynamic} frame}. That is the TLS record machinery -- handshake transcript,
record buffers, certificate-chain working storage -- laid out on the environment
task's stack, and the \texttt{dynamic} tag is the tool telling you the truth: part
of that frame is sized at run time, so the 5\,008-byte worst-case is a \emph{floor},
not a ceiling, and the \texttt{+ext} tag adds the (uncounted) runtime frames beneath
it. This is precisely the case the static pass cannot close on its own -- and
exactly why the environment-task stack for this example is set to 32\,KiB, a
deliberately generous margin over the $\sim$5\,KiB static floor precisely because
that floor is not the whole story, and why you confirm it with the painted
watermark rather than trusting the static sum. Run it under
\texttt{./x stack tls\_weather -{}-run} and the measured \texttt{stack:} line tells
you how much of that allotment a real fetch actually touched.

The footprint side, from \texttt{./x mem tls\_weather}:

\begin{lstlisting}
   flash (code + rodata + IRAM image):    287 519 B
   RAM   DRAM .data ............     23 468 B
   RAM   DRAM .bss / stacks .....    125 960 B
   RAM   total .................    153 268 B
\end{lstlisting}

\noindent
About 287\,KiB of flash (the TLS, X.509 and W5500 code is substantial) and roughly
153\,KiB of RAM, the bulk of it the $\sim$118\,KiB of \texttt{.bss} -- the static
session buffers, the W5500 socket storage and the heap arena, all reserved at link
time exactly as this chapter argued they should be. Nothing here grows on demand:
the whole budget, code and data and stack alike, is a fixed number you can read off
two commands and hold against the 512\,KiB of on-chip SRAM before the board ever
runs.
